\documentclass[12pt,a4paper]{article}
\usepackage[utf8]{inputenc}
\usepackage{amsmath,amssymb}
\usepackage{graphicx}
\usepackage[left=1.5cm,right=1cm,top=1.3cm,bottom=1.3cm]{geometry}
\usepackage{lastpage}
% Font Times New Roman (cần XeLaTeX + fontspec — uncomment nếu VPS hỗ trợ):
% \usepackage{fontspec}
% \setmainfont{Times New Roman}

\begin{document}

\begin{center}
\begin{tabular}{p{0.4\textwidth}p{0.6\textwidth}}
\textbf{1} & 5 \tabularnewline
\textbf{2} & \textbf{6} \tabularnewline
\textbf{3} & \textit{7, 8} \tabularnewline
\textit{4} & \tabularnewline
Họ và tên: ...................................................................................... \newline Số báo danh: .................................................................................. & \textbf{Mã đề: 1761} \tabularnewline
\end{tabular}
\end{center}
\vspace{0.5cm}

\begin{center}
{\Large\textbf{ĐÁP ÁN VÀ LỜI GIẢI}}
\end{center}

\noindent\textbf{PHẦN I. Câu trắc nghiệm nhiều phương án lựa chọn.}

\medskip

\noindent\textcolor{blue}{\textbf{Câu 1.}}

\noindent \textbf{Đáp án A.}

\noindent\textbf{Lời giải.}

Ta có $\displaystyle\int \left(\dfrac{1}{x} + \dfrac{1}{\cos ^2 x}\right) \mathrm{\,d}x = \displaystyle\int \dfrac{1}{x} \mathrm{\,d}x + \displaystyle\int \dfrac{1}{\cos ^2 x} \mathrm{\,d}x = \ln |x| + \tan x + C$.

\medskip
\noindent\textcolor{blue}{\textbf{Câu 2.}}

\noindent \textbf{Đáp án C.}

\noindent\textbf{Lời giải.}

Dựa vào đồ thị suy ra hàm số đồng biến trên khoảng $(-1;1)$.

\medskip
\noindent\textbf{PHẦN II. Câu trắc nghiệm đúng sai.}

\medskip

\noindent\textcolor{blue}{\textbf{Câu 1.}}

\noindent \textbf{Đáp án: ĐSĐĐ}

\noindent\textbf{Lời giải.}

Tập xác định $\mathscr{D}=\mathbb{R} \setminus \{2\}$.

		Bảng xét dấu
		
% [Hình TikZ không compile được: tikz_1d02414fd678]

		Vậy hàm số $y=f(x)$ nghịch biến trên các khoảng $(1;2)$ và $(2;3)$.

		Mà $f'(1)=0\Rightarrow 1-4-2b+2=0\Rightarrow b=-\dfrac{1}{2}$. Khi đó $f(x)=\dfrac{x^2-\dfrac{1}{2}x-2}{x-2}$. \\

		Mà hàm số $y=f(x)$ đồng biến trên đoạn $[3;4]$ nên $\max\limits_{[3;4]} f(x)=f(4)=6$.

\medskip
\noindent\textcolor{blue}{\textbf{Câu 2.}}

\noindent \textbf{Đáp án: SĐĐĐ}

\noindent\textbf{Lời giải.}

Biểu thức vận tốc tức thời của vật là $v(t)=-t^2+8t+9$.\\

		Vận tốc của vật tại các thời điểm $t=3$ giây là $v(3)=-3^2+8\cdot 3+9=24 \text{ (m/s)}$.\\

		Vật đứng yên khi
		
$$v(t)=0\Leftrightarrow -t^2+8t+9=0\Leftrightarrow \left[\begin{aligned}&t=-1 \,(\text{không thỏa mãn}) \\&t=9 \,(\text{thỏa mãn}).\end{aligned}\right.$$

		Trong khoảng thời gian từ lúc vật bắt đầu chuyển động đến khi vật đứng yên thì vật không đổi chiều chuyển động.\\

		Ta có $s(0)=0$ và $s(9)=-\dfrac{1}{3}\cdot 9^3+4\cdot 9^2+9\cdot 9=162 \text{ (m)}$.\\

		Quãng đường vật đi được từ lúc bắt đầu chuyển động đến thời điểm $t=9$ giây là $\left|s(9)-s(0)\right|=162 \text{ (m)}$.\\

		Biểu thức gia tốc tức thời của vật là $a(t)=v'(t)=-2t+8$.\\

		Gia tốc của vật tại thời điểm $t=3$ giây là $a(3)=-2\cdot 3+8=2\, \left(\text{m/s}^2\right)$.\\

		Phương trình $v'(t)=0$ có nghiệm $t=4$.\\

		Bảng biến thiên của $v(t)$ trên nửa khoảng $[0;+\infty)$\\

% [Hình TikZ không compile được: tikz_31955650d863]
Dựa vào bảng biến thiên trên, vật tăng tốc khi $t \in[0; 4]$.

\medskip
\noindent\textbf{PHẦN III. Câu trả lời ngắn.}

\medskip

\noindent\textcolor{blue}{\textbf{Câu 1.}}

\noindent \textbf{Đáp số: -1}

\noindent\textbf{Lời giải.}

Tập xác định $\mathscr{D}=[-1;1]$.\\

 Ta có $y'=\dfrac{-x}{\sqrt{1-x^2}}$ suy ra $y'>0$ khi $x<0$.\\

 Kết hợp với tập xác định, ta có khoảng đồng biến có độ dài lớn nhất là khoảng $(-1;0)$.\\

 Từ đó, ta có $a=-1$, $b=0$ suy ra $a+b=-1$.

\medskip
\noindent\textcolor{blue}{\textbf{Câu 2.}}

\noindent \textbf{Đáp số: -3}

\noindent\textbf{Lời giải.}

Hàm số xác định khi $x^2+3x>0\Leftrightarrow \left[\begin{array}{l}x>0\\x<-3.\end{array}\right.$\\

		Tập xác định $\mathscr{D}=(-\infty;-3)\cup(0;+\infty)$.\\

		Đạo hàm $y'=\dfrac{2x+3}{(x^2+3x)\ln 2}$.\\

		Hàm số nghịch biến khi $y'<0\Rightarrow 2x+3<0\Leftrightarrow x<-\dfrac{3}{2}$.\\

		Kết hợp điều kiện xác định, hàm số nghịch biến trên khoảng $(-\infty;-3)$.\\

		Vậy giá trị lớn nhất của $a$ là $-3$.

\medskip
\noindent\textbf{PHẦN IV. Câu tự luận.}

\medskip

\noindent\textcolor{blue}{\textbf{Câu 1.}}

\noindent\textbf{Lời giải.}

Vận tốc của vật là $v(t)=s'(t)=-3t^2+6t+6$.\\

		Ta có $v'(t)=-6t+6$.\\

		Cho $v'(t)=0\Leftrightarrow t=1$.\\

		Bảng biến thiên\\

% [Hình TikZ không compile được: tikz_bcb7f028b670]

		Dựa vào bảng biến thiên, ta thấy $v(t)$ tăng khi $t\in (0;1)$. Vậy vật tăng tốc trong khoảng thời gian từ gian $1$ giây tính từ lúc bắt đầu chuyển động.

\medskip
\noindent\textcolor{blue}{\textbf{Câu 2.}}

\noindent\textbf{Lời giải.}

Tốc độ truyền bệnh $y=f'(t)=-3t^2+90t$ với $t\in [0; 25]$.\\

		Ta có $y'=-6t+90$; $y'=0\Leftrightarrow t=15$.\\

		Bảng biến thiên
		
% [Hình TikZ không compile được: tikz_7d2da6ddaef5]

		Tốc độ truyền bệnh giảm từ ngày $15$ đến ngày $25$.\\

		Do đó $m=15$, $n=25$.\\
Suy ra $m-n=10$.

\medskip

\end{document}